\section{Authenticated updates: reissue and restructure}
\label{sec:authenticated-updates}

Replacing someone's token, or changing the shape of the split tree, has to
reach every store that held the old picture. Both changes ship as the same
sealed \file{.kqpb} envelope a private bridge uses, so the mailbox routes
them without being able to read them --- but the letter inside is
\emph{signed} by an authorizing label, so each store decides for itself
whether to apply it.

\subsection{The four-part authorization rule}

A store applies an update only when all four of these hold:

\begin{itemize}
  \item \textbf{Addressed here} --- the envelope's recipient key is an
  unrevoked encryption key this store has registered under the label named
  inside the letter. Re-addressing a letter to another mailbox gets it
  rejected.
  \item \textbf{Authorized} --- the signer is the subject itself or one of
  its ancestors in the dotted hierarchy (\code{M} over \code{M.S} over
  \code{M.S.2}), and this store already holds a registered signing key for
  that label. A peer branch cannot restructure your tree or swap your
  token.
  \item \textbf{Signed} --- Ed25519 over a domain-separated hash that covers
  every field, including the recipient label and its public key.
  \item \textbf{In order} --- a reissue must be exactly one past the last
  one this store applied for that person; a restructure must carry a public
  generation strictly greater than the one stored. Replays, stale
  envelopes, and skipped sequences change nothing.
\end{itemize}

\begin{lstlisting}
# The authorizing label needs a registered signing key:
keyquorum generate --type signing --label M --register \
  --public-key-out M.sign.pub > M.sign.key

# Restructure: after `add`, `revoke --evict`, or `bind`, announce the new
# shape. Each active leaf gets its own visible slice at the next generation.
# Root (M) is effective immediately. A label with a parent, such as M.S,
# writes a proposal and waits. --push uploads the envelopes in the same
# command, after the change is saved; each command stages them in its own
# folder under ./outbox. With --output-dir instead, give each command its own
# directory: these can overlap in which leaves they address (a proposal for
# M.S can name the same leaves the root's own restructure just did), and
# package writers never overwrite an existing file.
keyquorum tree restructure 1 --as M   --signing-key-file M.sign.key --push
keyquorum tree restructure 1 --as M.S --signing-key-file M.S.sign.key --push
keyquorum tree countersign 1 --as M   --signing-key-file M.sign.key --push
# the same countersignature, and a restructure, from a slot:
# keyquorum tree countersign 1 --as M --device ./usb --slot M --push
# keyquorum tree restructure 1 --as M --slot ./usb=M --push

# Reissue: M.S.2 lost their token and generated a replacement.
keyquorum reissue --node M.S.2 --key-id 1 \
  --encryption-public-key-file M.S.2.new.pub \
  --as M --signing-key-file M.sign.key --revoke-previous --push
\end{lstlisting}

\noindent Without \flag{--push}, give \flag{--output-dir} and upload the
folder with \code{relay push --dir}, which is still how envelopes that were
written offline are carried. \code{reissue} and \code{tree restructure} also
take \flag{--slot} to sign with a registered slot instead of a key file.

Each store applies these the same way it applies a bridge envelope --- the
kind byte in the header decides which it is:

\begin{lstlisting}
keyquorum --db M.S.1.sqlite inbox open
# Still possible, not recommended: download and import by hand with a key file
# keyquorum --db M.S.1.sqlite relay pull --import --share-file M.S.1.key
keyquorum --db M.S.1.sqlite bridge private import --file M.S.1.kqpb \
  --share-file M.S.1.key
keyquorum --db M.S.1.sqlite updates    # what this store has applied
\end{lstlisting}

\subsection{What a reissue changes}

A reissue reaches every store whose own slice named that person, plus every
party of every live private bridge they are on --- nobody else, since
nobody else held the key. It repoints that person's tree leaves and bridge
roster entries, and drops the leaf's sealed share, which was wrapped to the
retired token and cannot be opened by its replacement. Recover it from the
quorum, or reseal with \code{bind --public-key-file}.

\begin{cautionBox}
The shared secret of a private bridge is \emph{not} rotated by a reissue.
Run \code{bridge private remove-member} to roll that generation if the
retired token could have been compromised.
\end{cautionBox}

The subject's own copy is addressed to the \textbf{incoming} encryption
key, so their replacement device must have registered it (\code{generate
--register}, or \code{register}) before importing --- which is also what
lets it open the letter.

\subsection{Restructure vs.\ publish}

\code{tree restructure} differs from \code{tree publish} in exactly this
way: \code{publish} replaces a relay document with a store's own view,
while a restructure hands each person a slice they can verify the
authority for.

A restructure whose authorizer is the root is applied as soon as a store
accepts the envelope. A restructure whose authorizer has a parent
(\code{M.S} under \code{M}) is a proposal until that parent runs \code{tree
countersign}; import accepts that change only with the parent's
countersignature. A routine employee reissue by that employee's own parent
(\code{M.S} reissuing \code{M.S.1}) stays a single signature.
